
\documentclass[12pt]{article}
%%%%%%%%%%%%%%%%%%%%%%%%%%%%%%%%%%%%%%%%%%%%%%%%%%%%%%%%%%%%%%%%%%%%%%%%%%%%%%%%%%%%%%%%%%%%%%%%%%%%%%%%%%%%%%%%%%%%%%%%%%%%%%%%%%%%%%%%%%%%%%%%%%%%%%%%%%%%%%%%%%%%%%%%%%%%%%%%%%%%%%%%%%%%%%%%%%%%%%%%%%%%%%%%%%%%%%%%%%%%%%%%%%%%%%%%%%%%%%%%%%%%%%%%%%%%
\usepackage{amsmath}
\usepackage[compat2]{geometry}

\setcounter{MaxMatrixCols}{10}
%TCIDATA{TCIstyle=LaTeX article (bright).cst}

%TCIDATA{OutputFilter=LATEX.DLL}
%TCIDATA{Version=5.50.0.2953}
%TCIDATA{<META NAME="SaveForMode" CONTENT="1">}
%TCIDATA{BibliographyScheme=Manual}
%TCIDATA{Created=Tuesday, August 14, 2012 13:45:33}
%TCIDATA{LastRevised=Wednesday, September 15, 2021 12:44:13}
%TCIDATA{<META NAME="GraphicsSave" CONTENT="32">}
%TCIDATA{<META NAME="DocumentShell" CONTENT="Exams and Syllabi\SW\Assignment">}

\setlength{\topmargin}{-1.0in}
\setlength{\textheight}{9.25in}
\setlength{\oddsidemargin}{0.0in}
\setlength{\evensidemargin}{0.0in}
\setlength{\textwidth}{6.5in}
\def\labelenumi{\arabic{enumi}.}
\def\theenumi{\arabic{enumi}}
\def\labelenumii{(\alph{enumii})}
\def\theenumii{\alph{enumii}}
\def\p@enumii{\theenumi.}
\def\labelenumiii{\arabic{enumiii}.}
\def\theenumiii{\arabic{enumiii}}
\def\p@enumiii{(\theenumi)(\theenumii)}
\def\labelenumiv{\arabic{enumiv}.}
\def\theenumiv{\arabic{enumiv}}
\def\p@enumiv{\p@enumiii.\theenumiii}
\pagestyle{plain}
\setcounter{secnumdepth}{0}
\parindent=0pt
\input{tcilatex}
\begin{document}


\begin{center}
\begin{equation*}
\FRAME{itbpF}{1.3214in}{0.9496in}{0in}{}{}{Figure}{\special{language
"Scientific Word";type "GRAPHIC";maintain-aspect-ratio TRUE;display
"USEDEF";valid_file "T";width 1.3214in;height 0.9496in;depth
0in;original-width 1.3552in;original-height 0.9669in;cropleft "0";croptop
"1";cropright "1";cropbottom "0";tempfilename
'QYQBI30B.bmp';tempfile-properties "XPR";}}
\end{equation*}%
\textbf{Electromagnetismo - 230045}

\emph{Gu\'{\i}a N}$%
%TCIMACRO{\U{b0}}%
%BeginExpansion
{{}^\circ}%
%EndExpansion
2$\emph{\ - Campo el\'{e}ctrico producido por cargas puntuales}

\textit{Profesor: Arturo Fern\'{a}ndez P\'{e}rez}
\end{center}

\section{Problemas unidimensionales}

\begin{enumerate}
\item Dos cargas puntuales est\'{a}n situadas sobre el eje $X$. La primera
tiene una carga $+Q$ y est\'{a} en $x=-a$. La segunda, de carga desconocida,
est\'{a} situada en $x=3a$. El campo el\'{e}ctrico total en el origen tiene
magnitud total $2k_{e}Q/a^{2}$. \textquestiondown Qu\'{e} valores puede
tomar la carga desconocida? \textbf{R:} $-9Q.$

\item Para la configuraci\'{o}n de la Fig. 1, determine el punto (aparte del
infinito) donde el campo el\'{e}ctrico se anula. \textbf{R:} A $1,821$ [m] a
la izquierda de la carga $-2,5$ $\mu C.$%
\begin{equation*}
\FRAME{itbpFU}{2.4422in}{1.1026in}{0in}{\Qcb{Figura 1}}{}{Figure}{\special%
{language "Scientific Word";type "GRAPHIC";maintain-aspect-ratio
TRUE;display "USEDEF";valid_file "T";width 2.4422in;height 1.1026in;depth
0in;original-width 4.7911in;original-height 2.1473in;cropleft "0";croptop
"1";cropright "1";cropbottom "0";tempfilename
'QYQBI40C.wmf';tempfile-properties "XPR";}}
\end{equation*}

\item Determine la magnitud y direcci\'{o}n del campo el\'{e}ctrico que
equilibra la fuerza peso para el caso de:

\begin{itemize}
\item Un prot\'{o}n. \textbf{R}: $\vec{E}=1,025\times 10^{-7}$ $\hat{\jmath}$
[N/C]

\item Un electr\'{o}n. \textbf{R}: $\vec{E}=5,580\times 10^{-11}$ $\left( -%
\hat{\jmath}\right) $ [N/C]

\item Un ion de Li$^{+}$\pagebreak\ \textbf{R}: $\vec{E}=7,172\times 10^{-7}$
$\hat{\jmath}$ [N/C]
\end{itemize}
\end{enumerate}

\section{Problemas bidimensionales}

\begin{enumerate}
\item Tres cargas puntuales negativas se ubican como muestra la Fig. 2.\
Determine el campo el\'{e}ctrico resultante en el punto P. Luego, determine
la fuerza el\'{e}ctrica se ejerce sobre una carga negativa $q=-1$ [$\mu C]$
en el punto P, utilizando el resultado anterior. \textbf{R: }$\vec{E}%
=1.04\times 10^{7}$ $(-\hat{\imath})$ $[\frac{N}{C}]$ ; $\vec{F}=10.4$ $\hat{%
\imath}$ $[N]$%
\begin{equation*}
\FRAME{itbpFU}{1.8697in}{2.4656in}{0in}{\Qcb{Figura 2}}{}{Figure}{\special%
{language "Scientific Word";type "GRAPHIC";maintain-aspect-ratio
TRUE;display "USEDEF";valid_file "T";width 1.8697in;height 2.4656in;depth
0in;original-width 3.4134in;original-height 4.5143in;cropleft "0";croptop
"1";cropright "1";cropbottom "0";tempfilename
'QYQBI40F.wmf';tempfile-properties "XPR";}}
\end{equation*}

\item En la Fig. 3 se muestran 4 cargas puntuales, $q_{1}=q_{2}=+5q$, $%
q_{3}=+3q$ y $q_{4}=-12q$ donde $q=5$ [$\mu C$]. Si $d=5$ [mm]$,$ determine
el campo el\'{e}ctrico resultante en el punto P$.$ \textbf{R}: $\vec{E}=\vec{%
0}$ 
\begin{equation*}
\FRAME{itbpFU}{2.431in}{2.4344in}{0in}{\Qcb{Figura 3}}{}{Figure}{\special%
{language "Scientific Word";type "GRAPHIC";maintain-aspect-ratio
TRUE;display "USEDEF";valid_file "T";width 2.431in;height 2.4344in;depth
0in;original-width 3.8311in;original-height 3.8372in;cropleft "0";croptop
"1";cropright "1";cropbottom "0";tempfilename
'QYQBI40D.wmf';tempfile-properties "XPR";}}
\end{equation*}

\item La Fig. 4 muestra un arreglo de protones ($p$) y electrones ($e$) en
un arco circular de radio $R=2$ [cm], con \'{a}ngulos $\theta _{1}=30%
%TCIMACRO{\U{b0}}%
%BeginExpansion
{{}^\circ}%
%EndExpansion
$, $\theta _{2}=50%
%TCIMACRO{\U{b0}}%
%BeginExpansion
{{}^\circ}%
%EndExpansion
$, $\theta _{3}=30%
%TCIMACRO{\U{b0}}%
%BeginExpansion
{{}^\circ}%
%EndExpansion
$ y $\theta _{4}=20%
%TCIMACRO{\U{b0}}%
%BeginExpansion
{{}^\circ}%
%EndExpansion
$. Determine la magnitud y direcci\'{o}n del campo el\'{e}ctrico resultante
en el centro del arco. \textbf{R:} $\vec{E}=9,252\times 10^{-7}\hat{\imath}%
-3,820\times 10^{-6}\hat{\jmath}$ [N/C]%
\begin{equation*}
\FRAME{itbpFU}{2.4509in}{1.9225in}{0in}{\Qcb{Figura 4}}{}{Figure}{\special%
{language "Scientific Word";type "GRAPHIC";maintain-aspect-ratio
TRUE;display "USEDEF";valid_file "T";width 2.4509in;height 1.9225in;depth
0in;original-width 5.4587in;original-height 4.2739in;cropleft "0";croptop
"1";cropright "1";cropbottom "0";tempfilename
'QYQBI40E.wmf';tempfile-properties "XPR";}}
\end{equation*}

\item En la configuraci\'{o}n de la Fig. 5, considere $a=10$
[cm],\thinspace\ $b=15$ [cm] y $q=8,0$ [nC]$.$

\begin{enumerate}
\item Determine el campo el\'{e}ctrico resultante inducido por las cargas en
el punto P. \textbf{R: }$\vec{E}=-503,7\hat{\imath}-2469,14\hat{\jmath}$
[N/C]

\item Si una carga puntual $q_{2}=-30$ [nC] se ubica en dicho punto,
determine la fuerza el\'{e}ctrica neta ejercida sobre ella. \textbf{R: }$%
\vec{F}=15,1\hat{\imath}+74,1\hat{\jmath}$ [$\mu N$]%
\begin{equation*}
\FRAME{itbpFU}{2.3393in}{1.6483in}{0in}{\Qcb{Figura 5}}{}{Figure}{\special%
{language "Scientific Word";type "GRAPHIC";maintain-aspect-ratio
TRUE;display "USEDEF";valid_file "T";width 2.3393in;height 1.6483in;depth
0in;original-width 4.0811in;original-height 2.8669in;cropleft "0";croptop
"1";cropright "1";cropbottom "0";tempfilename '../../clases II-2020/230045
Electromagnetismo/QYQC4Q0G.wmf';tempfile-properties "XPR";}}
\end{equation*}
\end{enumerate}
\end{enumerate}

\end{document}
